\section{Introduction}

% - Mise en contexte (essor des IAG)
% - Problématique générale
% - Intérêt épistémologique
% - Annonce du plan

Les intelligences artificielles g\'en\'eratives (IAG) ont fait apparition dans l'espace public et \'educatif avec une rapidit\'e sans pr\'ec\'edent. 
Depuis le lancement grand public de ChatGPT fin 2022, ces outils capables de produire des textes, des raisonnements et des synth\`eses de documents se sont r\'epandus massivement dans les pratiques \'etudiantes. 
L'histoire de l'IA en \'education est pourtant ancienne: autour des années 1950, des premières expérimentations d'utilisation de l'IA dans le cadre scolaire ont \'et\'e men\'ees, notamment pour vérifier la justesse d'expressions math\'ematiques~\citep{siham2024intelligence}.

Dans les années 1980, l'éducation scolaire est devenue le terrain de recherche principal de l'IA, avec la création de systèmes modélisant les connaissances d'experts dans divers domaines.

Ce que les IA g\'en\'eratives introduisent de radicalement nouveau, c'est la capacit\'e de \emph{g\'en\'erer} du contenu en langage naturel qui ressemble \`a une production humaine: une dissertation, un rapport, un m\'emoire.

Cette \'evolution souligne une question fondamentale pour le syst\`eme \'educatif: si un outil peut produire en quelques secondes un travail formellement irr\'eprochable, que mesurent encore nos \'evaluations?
Et par cons\'equent, que signifie encore \emph{r\'eussir} \`a l'\'ecole, ou \emph{\^etre comp\'etent}~? Ce n'est pas l\`a une question purement pratique ou organisationnelle. 
Elle touche \`a des fondements \'epist\'emologiques: la nature du savoir scolaire, les crit\`eres de sa validation, la relation entre l'individu et les outils qu'il mobilise pour produire de la connaissance.

L'int\'er\^et \'epist\'emologique du sujet r\'eside pr\'ecis\'ement dans ce que les IAG font appara\^\i tre par contraste: des pr\'esuppos\'es sur l'apprentissage et l'\'evaluation que le syst\`eme scolaire avait rarement eu besoin d'expliciter.
L'apparition de l'IAG agit comme un r\'ev\'elateur, mettant en lumi\`ere des tensions que le syst\`eme \'educatif n'avait pas encore eu \`a affronter.
Ainsi, la probl\'ematique de ce m\'emoire n'est pas de proposer un syst\`eme \'educatif alternatif, mais de d\'ecrire et d'analyser ces tensions.

Notre plan sera le suivant. 
Nous commencerons par d\'efinir les notions centrales: comp\'etence, r\'eussite scolaire, et intelligences artificielles g\'en\'eratives (section~\ref{sec:definitions}).  % chktex 44
Nous poserons ensuite le cadre th\'eorique et \'epist\'emologique qui sous-tend notre analyse (\ref{sec:cadre}). 
Nous examinerons comment les IAG remettent en question les m\'ethodes d'\'evaluation actuelles (\ref{sec:remise}), avant d'exposer les deux grandes positions institutionnelles --- interdiction et int\'egration --- face \`a ce d\'efi (\ref{sec:positions}). 
Enfin, nous d\'egagerons les tensions \'epist\'emologiques profondes que ce ph\'enom\`ene r\'ev\`ele (\ref{sec:tensions}), avant de conclure par une discussion ouverte (\ref{sec:discussion}).
